\documentclass[../main.tex]{subfiles}
\begin{document}
\section{Visão da Empresa}

\subsection{Origem e Propósito}

A \textbf{CSIS} nasceu da convergência de duas realidades duras experimentadas diretamente por seu fundador na intersecção entre o mercado educacional e o corporativo. A primeira delas é a dor clássica do estudante de cibersegurança: o abismo prático entre acumular conhecimento técnico e de fato conseguir gerar renda, obter o primeiro emprego ou angariar contratos reais na área. A segunda é a observação crônica de que escritórios jurídicos e empresas possuem uma alta demanda por serviços forenses e de segurança, mas esbarram na lentidão e nos custos excessivos de consultorias tradicionais.

Para solucionar essa equação, a concepção original afastou-se do modelo burocrático de uma agência convencional. O propósito da CSIS não é restringir a equipe a um \textit{solopreneur} permanente, mas sim alicerçar uma estrutura infinitamente elástica e com baixa burocracia. O objetivo é permitir que novos talentos em ascensão possam ser integrados facilmente à rede de fluxo de trabalho livre (\textit{freelancing} com revisão), ganhando a chance de atuar em perícias reais e de ser remunerados financeiramente (modelo de \textit{Bounties}) enquanto a empresa escala sua entrega técnica e seu alcance comercial B2B.

\subsection{Missão, Visão e Valores}

\textbf{Missão}: Capacitar profissionais de todos os níveis a gerar renda real com cibersegurança através do aprendizado prático, enquanto atua provendo proteção, liderança técnica e inteligência forense acessível e de excelência para o mercado corporativo. \\

\textbf{Visão}: Consolidar-se, em 3 anos, como a maior comunidade educacional prática do país e a principal e mais confiável fonte de conteúdos estratégicos da área. Simultaneamente, firmar-se como a parceira de referência incontestável em perícia forense e soluções avançadas de mitigação de vulnerabilidades para os mercados jurídico e empresarial. \\

\textbf{Valores}: A atuação da CSIS apoia-se em princípios fundamentais e inegociáveis:
\begin{itemize}
    \item \textbf{Crescimento Colaborativo e Empatia}: O ecossistema se sustenta em acolher o iniciante, ensinar ativamente e incentivar o desenvolvimento mútuo;
    \item \textbf{Meritocracia Técnica}: Superação da teoria engessada pelas provas de conceito e resolução prática de \textit{Bounties}; o reconhecimento baseia-se unicamente na capacidade de entrega e no mérito;
    \item \textbf{Pragmatismo}: Foco irrestrito no que funciona sob as condições de estresse do mundo real, privilegiando a execução em cenários verdadeiros (\textit{Hands-on});
    \item \textbf{Sigilo e Ética Profissional}: Proteção estrita aos dados sensíveis corporativos e manutenção de elevados padrões de integridade moral e jurídica.
\end{itemize}


\subsection{Experiência e Capacidade Técnica}

A espinha dorsal de excelência da CSIS é garantida pelo \textit{background} técnico e prático de seu fundador, Christian Amarildo. Com um perfil singular que une a ponta ofensiva da cibersegurança à formalidade jurídica da investigação, Christian atua como Perito Judicial / Auxiliar da Justiça (TJPA) e traz expressiva experiência do setor de Perícia Criminal e Dev Blockchain na Polícia Civil do Estado do Pará. Em seu histórico prático, constam a elaboração de aproximadamente 450 laudos técnicos periciais, a construção de soluções em redes Blockchain para rastreabilidade de cadeias de custódia de evidências (Hyperledger Fabric) e atuações de peso logístico, como monitoramento e resposta a incidentes (\textit{Threat Hunting}) em centros de operação (SOC).

Todo o aparato prático está amparado por sólido arcabouço intelectual: Bacharelado em Ciência da Computação (UFPA) e MBAs super-intensivos em Computação Forense e Segurança da Informação, além de Cibersegurança e Gestão de Riscos (IPOG e Anhanguera). Essa estrutura é chancelada com certificações mundiais e capacitações corporativas da \textit{CompTIA Security+ ce}, Fortinet, Microsoft e AWS. Além das premiações em Hackathons de inovação aberta (Google Gemini), este arsenal multidisciplinar (que cruza habilidades de \textit{Red Team} como SQLMap, Metasploit, Burp Suite, com ferramentas de nível militar de \textit{Blue Team} como Cellebrite e Volatility) consagra a segurança técnica que chancelará a revisão superior (\textit{QA}) de cada entrega B2B da CSIS.

\subsection{Modelo de Atuação}

A atuação da organização foge radicalmente da consultoria tradicional operando num "Modelo de Rede e Esteira Forense Tática". Sem escritórios físicos ou pesadas folhas de pagamento de especialistas inativos, o cerne administrativo figura unicamente no fundador como \textit{Arquiteto de Soluções e Auditor Sênior de Qualidade}, coordenando a esteira que une a empresa e os alunos num laço contínuo:

\begin{enumerate}
    \item \textbf{Captação e Mapeamento (A Dor do Cliente):} A demanda técnica inicial chega através do mercado B2B, exigindo \textit{pentest} ou coleta/análise investigativa extrema. O escopo é fatiado sistemicamente;
    \item \textbf{Terceirização Prática Gamificada (O Bounty):} Fragmentado, o problema vai para uma arquitetura fechada no GitHub em formato de "\textit{Issue}" / Missão com prêmio associado. Via comunicação pelo Discord, o círculo restrito de alunos talentosos ("\textit{freelancers}") puxa a demanda, desenvolve métodos teóricos em laboratórios práticos e devolve uma documentação (\textit{Proof of Concept}) que responde com precisão à lacuna;
    \item \textbf{Entrega Controlada e Extração de Valor:} Por meio de um rígido \textit{Pull Request}, o fundador audita tudo. Se o arquivo estiver "padrão-ouro", o aluno é pago e o documento forense ganha a assinatura do Especialista e é despachado para a corporação judicial/cliente, resolvendo o problema original.
    \item \textbf{O Fechamento do Volante:} Todas essas técnicas inovadoras executadas, agora completamente esterilizadas de dados reais e anonimizadas, convertem-se em potentes tutorais técnicos \textit{Deep-Dive} publicados no YouTube e redes sociais para converter novos clientes pelo atestado vivo de autoridade corporativa.
\end{enumerate}
\end{document}